\scanpage{8}
Ceci posé, je dis que l'hom précédent
\[
\mathbf Rp_*(W)\longrightarrow e_*(\mathcal O_S)\otimes\tau_{B^*}[-n]
\]
est un isomorphisme. Il suffit pour ceci de\\
prouver que l'homomorphisme induit \deleted{\unsure{sur les}}\\
en tout point $\unsure{x|b}\in B^*$ par $\mathbf Li_{\unsure{x|b}}^*$ est un\\
isomorphisme (histoire de mapping-cylindres\ldots).\\
En les points $b\in B^*-\unsure{e_1|e}(S)$, ce n'est autre que le\\
lemme 2. En un point \deleted{$b\in\unsure{e}$} $b$ de la forme\\
$e(s)$, on est ramené au cas où $S$ est le spectre d'un\\
corps algébriquement clos, à faire une vérification de\\
compatibilité.....

\uline{Remarques} 1) a) $\Phi_{\mathcal L}$ transforme $e_A$ en $1_B$, $e_s$\\
($s$ section de $A/S$) en $\mathcal L_s$ (faisceau inv. alg. $\sim0$\\
sur $B$ rel $/S$), \inserted{donc} $e_s*F$ en $\mathcal L_s\otimes\Phi_{\mathcal L}(F)$.

b) $\Phi_{\mathcal L}$ transforme \deleted{$\omega_A[n]$ en $\unsure{e|\varepsilon}$, $\omega_A[n]$}\\
$\mathcal L_t$ ($t$ section de $B$ sur $S$) en $\tau[-n]e_{-t}$ i.e.\\
$\mathcal L_t\omega[n]$ en $e_{-t}$, \deleted{donc} $F\otimes\mathcal L_t\deleted{\unsure{\omega[n]}}$ en $\{\Phi_{\mathcal L}(F)*e_t,\deleted{\unsure{\omega[n]e}}$.

c) $\varepsilon_B(\Phi_{\mathcal L}(F))=\chi_A(F)$, $\chi_B(\Phi_{\mathcal L}(F))=(-1)^n\varepsilon_A(F)$.
\margin{$\varepsilon=$ \unsure{r-g|rang} virtuel\\
égalité dans \deleted{\unsure{le}}\\
$H^0(S,\unsure{\mathbb Z|\mathbf Z})$}

2) $\det^*\circ\Phi_{\mathcal L}$ définit \deleted{\unsure{sur les VA relatives}}
\[
\varphi:\underline{\operatorname{Pic}}_{A/S}\longrightarrow\underline{\operatorname{Pic}}_{B/S}
\]
qui, pour toute section $\delta$ de $\underline{\operatorname{NS}}_{A/S}$, définit
\[
\underline{\operatorname{Pic}}^\delta_{A/S}\longrightarrow\underline{\operatorname{Pic}}^{-\delta'}_{B/S}
\]
compatible avec
\margin{$-(-\delta')=\delta'$}
\mspage{7}
